% Thermodynamic Consistency of the Solubility and Vapor Pressure of a Binary
% Saturated Salt + Water System. 1. LiCl + H2O
% Dewen Zeng and Jun Zhou
% J. Chem. Eng. Data 2006, 51, 315-321   DOI: 10.1021/je050322o
%
% LaTeX transcription. Content and numerical values reproduced as printed.

\documentclass[11pt]{article}

\usepackage[T1]{fontenc}
\usepackage[utf8]{inputenc}
\usepackage{amsmath}
\usepackage{amssymb}
\usepackage{graphicx}
\usepackage{booktabs}
\usepackage{array}
\usepackage[margin=1in]{geometry}
\usepackage{authblk}
\usepackage[numbers]{natbib}

% superscript reference style, as in the original journal article
\newcommand{\rc}[1]{\textsuperscript{\citenum{#1}}}

\title{Thermodynamic Consistency of the Solubility and Vapor Pressure of a
Binary Saturated Salt + Water System. 1.\ LiCl + H$_2$O}

\author[1]{Dewen Zeng\thanks{To whom correspondence should be addressed.
E-mail: dewen\_zeng@hotmail.com. Fax: +86-731-8713642.}}
\author[2]{Jun Zhou}
\affil[1]{College of Chemistry and Chemical Engineering, Hunan University,
410082 Changsha, PR China}
\affil[2]{Central-South University, 410083 Changsha, People's Republic of China}

\date{Received for review August 13, 2005. Accepted December 11, 2005.}

\begin{document}

\maketitle

\begin{abstract}
\noindent
Solubility data of LiCl$\cdot n$H$_2$O ($n = 0$, 1, 2) was evaluated by
checking the consistency of evaluated vapor pressures with the experimental
data. This process includes the following: (1) A series of vapor pressure data
above the LiCl$\cdot n$H$_2$O saturated solution were evaluated, and reliable
data were selected as ``standard'' data. (2) A BET model was selected to fit
the selected experimental data of water activity at salt concentrations below
20 mol$\cdot$kg$^{-1}$ and at temperatures ranging from 273.15 K to 428.65 K.
(3) Vapor pressures for the solubility data given by different authors were
calculated with the BET model and further compared with the standard data.
(4) Those solubility data at which vapor pressures calculated with the BET
model obviously deviate from the standard data were considered to be
unreliable.
\end{abstract}

\section*{Introduction}

A large number of experimental solubility data for the binary system
LiCl + H$_2$O have been determined so far.\rc{ref1}\textsuperscript{--}\rc{ref4}
However, experimental data reported by different authors at different times
usually do not agree with each other, and the differences are quite large at
certain temperatures. To determine a reliable set of solubility values of this
system from a large number of experimental data, researchers have used
different evaluation methods. Link and Seidell\rc{ref1} took the average
values of many experimental data; however, Monnin et al.\rc{ref5} developed an
evaluation approach by plotting experimental points in
composition--temperature diagrams, where values outside the general trend were
rejected. Cohen-Adad\rc{ref4} compiled the solubility data of this system, but
as pointed out by Monnin et al., the criteria that Cohen-Adad retained for the
data selection are not clear.

However, the vapor pressures of the LiCl$\cdot n$H$_2$O saturated solution
have been measured by many
chemists,\rc{ref6}\textsuperscript{--}\rc{ref10} and more reliable
data\rc{ref11}\textsuperscript{--}\rc{ref16} have been reported.
Thermodynamically, if an experimental date of solubility (composition and
temperature) is correct, then the vapor pressure of the corresponding
saturated solution (composition and temperature) should also be correct. As
long as a thermodynamic model can accurately describe the thermodynamic
properties of a system, including water activity and solubility, one can
calculate the water activity at a given solubility point and therefore the
vapor pressure. Furthermore, by comparing the reliable vapor pressures with
the calculated ones for the given experimental solubility data, one could
evaluate the correctness of experimental solubility data.

Among the thermodynamic models\rc{ref17,ref19,ref20,ref21,ref22,ref23} for
concentrated electrolyte solution, we select the Pitzer--Simonson--Clegg
model\rc{ref17}\textsuperscript{--}\rc{ref19} and the modified BET
model\rc{ref20} to describe and predict the water activity of the
LiCl + H$_2$O system. Emphasis will be placed on the check of the prediction
ability of the models of the salt concentration range near a saturated
solution. Then, the vapor pressure data above the LiCl$\cdot n$H$_2$O
saturated solution will be compiled and evaluated; among them, reliable data
will be selected as ``standard'' data. After that, the vapor pressures of the
LiCl$\cdot n$H$_2$O saturated solution at given solubility points (components
and temperatures) will be calculated and compared with the standard data.
Those solubility data at which the calculated vapor pressures deviate
obviously from the standard data will be considered to be unreliable. In this
approach, the reliability of a series of solubility data will be evaluated.

\section*{Models for Concentrated Electrolyte Solution}

\subsection*{Pitzer--Simonson--Clegg Model}

The Pitzer--Simonson--Clegg model\rc{ref17}\textsuperscript{--}\rc{ref19}
expresses the excess Gibbs energy of an aqueous electrolyte system with a
long-range electrostatic term and a short-range Margules expansion, as
described by eq 14 in ref \citenum{ref19}.

For a binary MX + H$_2$O system, they derived the expression of the water
activity coefficient $f_1$ as

% Note: in the source PDF the last term reads "x_1 x_I 2 (2 - 3x_1) V_{1,MX}";
% the isolated "2" appears to be an OCR artifact and is omitted here.
\begin{equation}
\ln f_1 = \frac{2 A_x I_x^{3/2}}{1 + F I_x^{1/2}}
- x_{\mathrm{M}} x_{\mathrm{X}} B_{\mathrm{MX}} \exp\!\left(-\alpha I_x^{1/2}\right)
+ x_{\mathrm{I}}^{2}\left( W_{1,\mathrm{MX}}
+ (x_{\mathrm{I}} - x_1) U_{1,\mathrm{MX}}
+ x_1 x_{\mathrm{I}} (2 - 3x_1) V_{1,\mathrm{MX}} \right)
\end{equation}

where $F = 2150(d_1/DT)^{1/2}$; $x_{\mathrm{I}} = x_{\mathrm{M}} +
x_{\mathrm{X}} = 1 - x_1$; $A_x$ and $I_x$ are the Debye--H\"uckel parameter
and ionic strength based on the mole fraction; $x_a$, $x_c$, $d_1$, $D$, and
$T$ are the mole fractions of the anion and cation, the density of the solvent
water, the dielectric constant of the solvent, and the thermodynamic
temperature, respectively. $B_{\mathrm{MX}}$, $W_{1,\mathrm{MX}}$,
$U_{\mathrm{U,MX}}$ and $V_{1,\mathrm{MX}}$ are model parameters.

Setting $\alpha = 13$, we fit the model parameters in eq 1 to the experimental
water activity\rc{ref24} at 373.15 K and at LiCl concentrations below 18.585
mol$\cdot$kg$^{-1}$ and obtained the model parameter values in Table 1. The
dashed line in Figure 1 shows that the Clegg--Pitzer model can represent the
water activity at LiCl concentrations below 18.58 mol$\cdot$kg$^{-1}$ very
well.

\begin{table}[htbp]
\centering
\caption{Clegg--Pitzer Model Parameters of the System LiCl + H$_2$O at
373.15 K}
\begin{tabular}{lcccccc}
\toprule
solute & \begin{tabular}[c]{@{}c@{}}maximal\\ molality\end{tabular} & $\alpha$
& $B_{\mathrm{MX}}$ & $W_{1,\mathrm{MX}}$ & $U_{1,\mathrm{MX}}$
& $V_{1,\mathrm{MX}}$ \\
\midrule
LiCl & 18.585 & 13 & $-47.092$ & $-2.083$ & 14.5 & $-15.7479$ \\
\bottomrule
\end{tabular}
\end{table}

With the parameters in Table 1, the water activity is calculated up to the
LiCl saturated solution, as shown by the dashed line in Figure 1. Until now,
no water activity of the LiCl saturated solution has been directly measured at
373.15 K; however, the vapor pressures and solubility of the LiCl saturated
solution were measured by Applebey et al.\rc{ref11} over a wide temperature
range. Interpolating their data yields the vapor pressure of the LiCl
saturated solution at 373.15 K (10.2 kPa) and the solubility data
(30.39 mol$\cdot$kg$^{-1}$). Furthermore, the vapor pressure $p$ is converted
to water activity $a_{\mathrm{w}}$ according to eq
2\rc{ref27,ref28,ref29}

\begin{equation}
\ln a_{\mathrm{w}} = \ln\!\left( \frac{p}{p_0} \right)
+ \frac{B_T (P - P_0)}{RT}
\end{equation}

where $p_0$ is the vapor pressure of pure water at temperature $T$,
$R = 8.314$ Pa$\cdot$m$^3\cdot$mol$^{-1}\cdot$K$^{-1}$, and $B_T$ is the
second virial coefficient taken from ref \citenum{ref30} at different
temperatures. The analyzed water activity and solubility data are presented in
Figure 1. Greenspan reported that the relative humidity ($P/P_0$) over LiCl
saturated solution is $(9.90 \pm 0.77)$ \% at 373.15 K, from which the
analyzed $a_{\mathrm{w}}$ with eq 2 and the solubility
data\rc{ref1,ref3,ref4,ref5} are also presented in Figure 1.

One can see in Figure 1 that the predicted water activity at the saturated
point with the Pitzer--Simonson--Clegg model parameters deviates from various
analyzed data by about 0.01. That means that the prediction ability of the
model is good but its prediction accuracy is not high enough to give a proper
judgment of the correctness of different experimental data.

\begin{figure}[htbp]
\centering
% \includegraphics[width=0.6\textwidth]{figure1}
\caption{Water activity of the system LiCl + H$_2$O at 373.15 K:
- - -, results recalculated with the Pitzer--Simonson--Clegg model parameters
in Table 1; ---, results recalculated with the BET model; $\circ$,
$T = 373.15$ K, ref \citenum{ref23}. Other symbols, water activity converted
from Greenspan\rc{ref16} and solubility data from the literature:
$\triangle$, ref \citenum{ref1}; right-facing triangle, ref \citenum{ref3};
$\triangledown$, ref \citenum{ref4}; left-facing triangle,
ref \citenum{ref5}; I, uncertainty; $\bullet$, data interpolated from
Applebey et al.\rc{ref11}}
\end{figure}

\subsection*{BET Model}

Again, we fit the BET model\rc{ref20} to the literature data of water activity
at 373.15 K at salt concentration below 18.585 mol$\cdot$kg$^{-1}$. As shown
in Figure 2, the linear relation of the BET equation (eq 3) holds at water
activities below 0.4,

\begin{equation}
\frac{a_{\mathrm{w}} m}{55.51(1 - a_{\mathrm{w}})}
= \frac{1}{cr} + \frac{(c - 1) a_{\mathrm{w}}}{cr}
\end{equation}

where $a_{\mathrm{w}}$ and $m$ are the water activity and salt concentration
in mol$\cdot$kg$^{-1}$ and $r$ and $c$ are model parameters. The calculation
of the slope $(c - 1)/cr$ and intercept $1/cr$ of the line yields the
parameter values $r = 3.278$ and $c = 9.329$. With the BET parameters, we
predict the water activity of the system up to the saturated solution (solid
line in Figure 1). The predicted water activity of LiCl saturated solution
deviates from different experimental data by less than 0.003. Thus, the BET
model seems to have a stronger prediction ability than the former in very
concentrated solution. Consequently, in the following work, we will select the
BET model to simulate the properties of the LiCl + H$_2$O system in
unsaturated solution and evaluate its properties in saturated solution.

\begin{figure}[htbp]
\centering
% \includegraphics[width=0.6\textwidth]{figure2}
\caption{BET relation in the system LiCl + H$_2$O at 373.15 K: $\circ$,
converted values from the literature data;\rc{ref24} ---, fitted line at water
activities below 0.4.}
\end{figure}

With the same method shown in Figure 2, we fitted the BET equation (eq 3) to
the literature data\rc{ref24} of water activity at 273.15 K, 298.15 K,
323.15 K, and 373.15 K and the literature data\rc{ref31} at 428.65 K and
obtained BET parameters $r$ and $\Delta E$ as a function of temperature $T$
(Figure 3, Figure 4, and Figure 5): $r = 4.7323 - 0.00378(T/\mathrm{K})$,
$\Delta E/\mathrm{J}\cdot\mathrm{mol}^{-1} = R(T/\mathrm{K}) \ln c =
-8166.6 + 3.526(T/\mathrm{K})$. With increasing temperature, the maximal
hydration sites get slightly smaller, as well as the absolute value of the
energy change $|\Delta E|$ accompanying the movement of 1 mol of water
molecules from pure water onto the hydration sites of the LiCl salt. The
predicted water activities with the above BET parameters are presented in
Figure 6 and agree quite well with literature
data\rc{ref3,ref4,ref9,ref24,ref31} in concentrated solution over a wide
temperature range from 323.15 K to 428.65 K.

\begin{figure}[htbp]
\centering
% \includegraphics[width=0.6\textwidth]{figure3}
\caption{BET relation in the system LiCl + H$_2$O: ---, fitted to the
reference data; $\square$, $+$, $\triangle$, $\diamond$, $\triangledown$,
ref \citenum{ref24}; $\circ$, ref \citenum{ref31}.}
\end{figure}

If we assume that the BET linear relation (eq 3) holds over the whole
concentration range of LiCl solution up to its saturated solution, then the
prediction accuracy of the model should depend on the accuracy of the water
activity data selected for the parameter fitting. Gibbard and Scatchard's
data\rc{ref24} deviate markedly from that of Pearce and Nelson;\rc{ref7} the
latter were obtained by a dynamic method with lower accuracy. However, Gibbard
and Scatchard's data\rc{ref24} agree very well with Robinson and Stokes'
data\rc{ref26} at 298.15 K and differ from Kangro and Groeneveld's osmotic
coefficient\rc{ref25} at 298.15 K by only $\pm0.02$. The osmotic coefficients
determined by Kangro and Groeneveld, according to their statement, have an
accuracy of $\pm0.001$, thus it is reasonable to believe that the osmotic
coefficient data reported by Gibbard and Scatchard\rc{ref24} have an accuracy
of $\pm0.02$. This accuracy can be interpreted as water activity equal to
$0.130 \pm 0.002$ for a 18.542 mol$\cdot$kg$^{-1}$ LiCl solution at 298.15 K
or as the mass percent LiCl solubility $(40.8 \pm 0.1)$ \% at 298.15 K. With
such accuracy, the BET model can be used to judge the reliability of
solubility data from different sources; for example, the mass percent LiCl
solubility reported by Woskresenskaya and Yanatieva\rc{ref39} is 38.7 \% at
273.15 K, which is 2 \% lower than the 40.87 \% data given by Friend and
Colley.\rc{ref37}

\begin{figure}[htbp]
\centering
% \includegraphics[width=0.6\textwidth]{figure4}
\caption{BET parameter $\Delta E$ vs temperature in the system
LiCl + H$_2$O: ---, fitted line; $\square$, $+$, $\triangle$, $\diamond$,
$\triangledown$, ref \citenum{ref24}; $\circ$, ref \citenum{ref31}.}
\end{figure}

\begin{figure}[htbp]
\centering
% \includegraphics[width=0.6\textwidth]{figure5}
\caption{BET parameter $r$ vs temperature in the system LiCl + H$_2$O:
---, fitted line; $\square$, $+$, $\triangle$, $\diamond$, $\triangledown$,
ref \citenum{ref24}; $\circ$, ref \citenum{ref31}.}
\end{figure}

\begin{figure}[htbp]
\centering
% \includegraphics[width=0.6\textwidth]{figure6}
\caption{Water activity vs salt concentration in the system LiCl + H$_2$O:
---, recalculated values with the BET model. Symbols represent literature
data: $\square$, $T = 273.15$ K, ref \citenum{ref24}; $+$, $T = 298.15$ K,
ref \citenum{ref24}; $\triangle$, $T = 323.15$ K, ref \citenum{ref24};
$\diamond$, $T = 348.15$ K ref \citenum{ref24}; $\triangledown$,
$T = 373.15$ K, ref \citenum{ref24}; $\circ$, $T = 428.65$ K,
ref \citenum{ref31}. The following data are not used for parameter fitting:
$\blacktriangle$, $T = 323.15$ K, ref \citenum{ref9}; $\blacktriangledown$,
$T = 373.15$ K, salt concentration taken from the solubility data\rc{ref4} and
water activity converted with eq 2 according to vapor pressure
data;\rc{ref11} $\bullet$, $T = 428.65$ K, salt concentration taken from the
solubility data\rc{ref3} and water activity converted with eq 2 according to
vapor pressure data.\rc{ref11}}
\end{figure}

\section*{Evaluation of Experimental Data of the Vapor Pressure above
Saturated Electrolyte Solutions}

Before evaluating different solubility data, one must first determine a
``correct'' set of vapor pressure data of the LiCl$\cdot$2H$_2$O saturated
solution. So far, many vapor pressures or relative humidity data of the
LiCl$\cdot$2H$_2$O saturated solution have been
reported\rc{ref6,ref8,ref10,ref12,ref15,ref16} (Figure 7). Among them, the
data given by H\"uttig and Reuscher\rc{ref6} at low temperatures are obviously
not accurate; for example, they obtained the same vapor pressure value of
0.10665 kPa above the LiCl$\cdot$2H$_2$O saturated solution at 273.15 K and
278.15 K and the same vapor pressure value of 0.2666 kPa at 285.65 K and
288.15 K. The absolute accuracy of the relative humidity data of the LiCl
saturated solution cited by Rockland\rc{ref8} is only 1 \%.
Greenspan\rc{ref16} accumulated experimental data from various researchers and
calculated ``best'' values of relative humidity, but his analyzed data differ
remarkably from the experimental data given by Hedlin\rc{ref12} and
Acheson\rc{ref13} (Figure 7). Hedlin et al.\rc{ref12} stated that the
experimental uncertainty in their relative humidity is only
$\pm0.1$ \%,\rc{ref12} and Acheson\rc{ref13} reported that their vapor
pressure data above the LiCl saturated solution is very accurate with an
absolute uncertainty of only 1 Pa. Thirty-four years later, Morillon
et al.\rc{ref15} determined the vapor pressures of the LiCl$\cdot$2H$_2$O
saturated solution again; these data have an uncertainty of 2 \% and agree
with the data of Hedlin\rc{ref12} quite well. Therefore, the experimental data
of Hedlin et al.,\rc{ref12} Acheson,\rc{ref13} and Morillon
et al.\rc{ref15} are believed to be more reliable in this work and are
consequently fitted with a temperature function of the vapor pressure above
the LiCl$\cdot$2H$_2$O saturated solution:

\begin{equation}
p/\mathrm{kPa} = \exp\!\left( a + b/(T/\mathrm{K}) + c \ln(T/\mathrm{K}) \right)
\end{equation}

Equation 4 was derived by assuming that the heat of vaporization of water in
the solution depends linearly on temperature.\rc{ref32} Parameter values $a$,
$b$, and $c$ are tabulated in Table 2.

\begin{figure}[htbp]
\centering
% \includegraphics[width=0.6\textwidth]{figure7}
\caption{Vapor pressures of LiCl$\cdot$2H$_2$O saturated solution vs
temperature: $\bullet$, ref \citenum{ref12}; $\blacktriangle$,
ref \citenum{ref13}; $\blacksquare$, ref \citenum{ref15}; $\square$,
ref \citenum{ref8}; $\circ$, ref \citenum{ref10}; $\triangle$,
ref \citenum{ref16}; I, uncertainty; ---, standard line fitted only to the
reference data.\rc{ref12,ref13,ref15}}
\end{figure}

\begin{table}[htbp]
\centering
\caption{Fitted Standard Curve of the Vapor Pressure of the
LiCl$\cdot n$H$_2$O Saturated Solution with the Temperature}
\footnotesize
\begin{tabular}{llcc}
\toprule
phase & vapor pressure of saturated solution, $p$/kPa
& \begin{tabular}[c]{@{}c@{}}temperature\\ range, $T$/K\end{tabular}
& \begin{tabular}[c]{@{}c@{}}source of\\ data, refs\end{tabular} \\
\midrule
LiCl$\cdot$2H$_2$O(s) & $\ln(p/\mathrm{kPa}) = 133.6116 - 9278.12/(T/\mathrm{K}) - 18.19053 \ln(T/\mathrm{K})$ & 273--293 & 12, 15, 16 \\
LiCl$\cdot$H$_2$O(s) & $\ln(p/\mathrm{kPa}) = 149.41106 - 11419.74114/(T/\mathrm{K}) - 19.68268 \ln(T/\mathrm{K})$ & 293--268 & 11, 16 \\
LiCl(s) & $\ln(p/\mathrm{kPa}) = 70.83197 - 8611.23808/(T/\mathrm{K}) - 7.66919 \ln(T/\mathrm{K})$ & 368--453 & 11 \\
\bottomrule
\end{tabular}
\end{table}

The vapor pressure data above or the water activity of the
LiCl$\cdot$H$_2$O saturated solution has been reported by different
researchers.\rc{ref6,ref7,ref8,ref9,ref11,ref12,ref16,ref33} The data
reported by H\"uttig et al.\rc{ref6} and Applebey et al.\rc{ref11} are
slightly scattered. As mentioned above, Pearce and Nelson's data\rc{ref7} are
less reliable. The data cited by Rockland\rc{ref8} are not very accurate. The
data from Hedlin et al.,\rc{ref12} Acheson,\rc{ref13} and
Greenspan\rc{ref16} agree with each other quite well in their stated
uncertainty ranges. These data, along with Applebey et al.'s data above
350 K,\rc{ref11} are selected to fit the standard curve with eq 4, as shown in
Figure 8.

Unfortunately, the data for the vapor pressure above the LiCl(s) saturated
solution is available only from Applebey et al.\rc{ref11} The uncertainty of
their data is unknown. These data are fitted with eq 4, as shown in Figure 9.

\begin{figure}[htbp]
\centering
% \includegraphics[width=0.6\textwidth]{figure8}
\caption{Vapor pressures of LiCl$\cdot$H$_2$O saturated solution vs
temperature: $\square$, ref \citenum{ref13}; $\triangle$,
ref \citenum{ref12}; $\circ$, ref \citenum{ref16}; $\diamond$,
ref \citenum{ref11}; I, uncertainty; literature data; ---, fitted standard
line.}
\end{figure}

\begin{figure}[htbp]
\centering
% \includegraphics[width=0.6\textwidth]{figure9}
\caption{Vapor pressures of LiCl saturated solution vs temperature: $\circ$,
ref \citenum{ref11}; ---, fitted standard line.}
\end{figure}

\section*{Comparison of Vapor Pressure and Evaluation of Solubility Data}

As introduced above, various solubility data of LiCl$\cdot n$H$_2$O have been
published. Linke and Seidell\rc{ref1} compiled all of the existing solubility
data before 1965 and gave one set of average values of the experimental data.
After that, Clynne and Potter\rc{ref3} measured the solubility data of this
system again, claiming that their data have an average precision of
$\pm0.09$ \%. Recently, Cohend-Adad\rc{ref4} and Monnin et al.\rc{ref5}
evaluated all published solubility data again and gave results different from
those given by Linke and Seidell\rc{ref1} and by Clynne and
Potter.\rc{ref3} To evaluate their solubility data and other older data, we
will calculate the water activity at each solubility data point with the BET
model and then convert the water activity data to vapor pressure data with
eq 2. The calculated vapor pressures will be compared with the standard curve.

\subsection*{On the Solubility Data of LiCl$\cdot$2H$_2$O}

The vapor pressures calculated for the typical solubility data of Linke and
Seidell\rc{ref1}, Clynne and Potter,\rc{ref3} and Monnin et al.\rc{ref5} are
presented in Figure 10. Also presented in this Figure are the standard curve
and the BET model uncertainty arising from that of the experimental
data,\rc{ref24} namely, $\pm0.02$ for the osmotic coefficient. One can see
that most of the vapor pressures calculated for the solubility
data\rc{ref1,ref3,ref5} fall in the estimated uncertainty range of the
standard line, except for two points at 278.15 K and 273.15 K.\rc{ref1}

Furthermore, the vapor pressures of the LiCl$\cdot$2H$_2$O saturated solution
are calculated for solubility data given by other authors, and the comparison
with the standard value is presented in Figure 11. It is surprising that most
of the evaluated vapor pressures for the solubility
data\rc{ref35,ref36,ref38,ref39,ref40,ref42} fall outside the uncertainty
range of the standard values. The vapor pressures calculated for the
solubility data\rc{ref38,ref40} are smaller than the standard values, whereas
those for other solubility data\rc{ref35,ref36,ref39,ref42} are larger. An
obvious positive deviation of the vapor pressure from the standard curve
appears for the solubility data.\rc{ref35,ref39,ref42} The International
Critical Tables\rc{ref42} in 1928 reported that the eutectic temperature of
LiCl$\cdot$2H$_2$O + LiCl$\cdot$H$_2$O lies at 285.65 K, at which the
evaluated vapor pressure deviates markedly from the standard value (Figure 11
and Table 3). At the invariant points reported by other
researchers\rc{ref1,ref5,ref34} the vapor pressures are also calculated. All
of the calculated data and their differences from the standard curve are
tabulated in Table 3.

\begin{figure}[htbp]
\centering
% \includegraphics[width=0.6\textwidth]{figure10}
\caption{Vapor pressures of the LiCl$\cdot$2H$_2$O saturated solution vs
temperature: ---, standard values; \ldots, model uncertainty based on the
standard values; $\circ$, ref \citenum{ref1}; $\square$, ref \citenum{ref3};
$\triangle$, ref \citenum{ref5}. All symbols are BET model values evaluated
for the solubility data of the reference literature.}
\end{figure}

\begin{figure}[htbp]
\centering
% \includegraphics[width=0.6\textwidth]{figure11}
\caption{Vapor pressures of LiCl$\cdot$2H$_2$O saturated solution vs
temperature: ---, standard values; \ldots, model uncertainty basing on the
standard values; $\circ$, ref \citenum{ref11}; $\square$ containing $\times$,
ref \citenum{ref35}; $+$, ref \citenum{ref36}; $\triangle$,
ref \citenum{ref37}; $\diamond$, ref \citenum{ref38}; left-facing triangle,
ref \citenum{ref39}; $\triangledown$, ref \citenum{ref40}; $\bullet$,
ref \citenum{ref41}; $\times$, ref \citenum{ref42}. All symbols are BET model
values evaluated for the solubility data of the reference literature.}
\end{figure}

\begin{table}[htbp]
\centering
\caption{Vapor Pressures Evaluated at the LiCl$\cdot$2H$_2$O +
LiCl$\cdot$H$_2$O Peritectic Point and Their Comparisons with the Standard
Values}
\begin{tabular}{cccccc}
\toprule
\multicolumn{2}{c}{literature data of solubility}
& \multicolumn{2}{c}{vapor pressure, kPa} & & \\
\cmidrule(r){1-2} \cmidrule(r){3-4}
$m_{\mathrm{LiCl}}$/mol$\cdot$kg$^{-1}$ & $T$/K & evaluated & standard
& deviation\textsuperscript{a}, kpa & \\
\midrule
19.57\textsuperscript{1}   & 291.65\textsuperscript{1}  & 0.233 & 0.236 & 0.003 & \\
19.42\textsuperscript{34}  & 292.55\textsuperscript{34} & 0.250 & 0.246 & $-0.003$ & \\
19.57\textsuperscript{5}   & 293.94\textsuperscript{5}  & 0.272 & 0.263 & 0.009 & \\
17.599\textsuperscript{42} & 285.65\textsuperscript{42} & 0.187 & 0.177 & $-0.011$ & \\
\bottomrule
\end{tabular}

\vspace{2pt}
\begin{minipage}{0.8\textwidth}
\footnotesize \textsuperscript{a} Deviation = standard value minus evaluated
value.
\end{minipage}
\end{table}

\subsection*{On the Solubility Data of LiCl$\cdot$H$_2$O}

The vapor pressures calculated for the LiCl$\cdot$H$_2$O solubility
data\rc{ref1,ref3,ref5,ref43} and their comparison with the standard curve are
presented in Figure 12. The calculated vapor pressures for all four sets of
solubility data agree with the standard values very well, with the average
percent deviation from the standard curve being smaller than 1.1 \% (Table 4).
It is noteworthy that although these authors used different methods to
evaluate and obtain their data they found consistent results. However, the
deviations of the vapor pressures calculated for other experimental solubility
data\rc{ref6,ref9,ref35,ref37,ref42} from the standard values are a little
larger (Figure 13 and Table 4), with the average percent deviation being
larger than 1.5 \%. The vapor pressures calculated for the most recently
finished two sets of experimental solubility data\rc{ref3,ref43} agree better
with the standard values than for the older experimental
data.\rc{ref6,ref9,ref35,ref37,ref42} On account of the uncertainty
($\pm0.02$) in the osmotic coefficient, to which the BET model parameters are
fitted, only those sets of solubility data with average deviations larger than
2 \% in Table 4 are believed to be unreliable.

\begin{figure}[htbp]
\centering
% \includegraphics[width=0.6\textwidth]{figure12}
\caption{Vapor pressures of the LiCl$\cdot$H$_2$O saturated solution vs
temperature: ---, standard values; \ldots, model uncertainty based on the
standard values; $\circ$, ref \citenum{ref1}; $\square$, ref \citenum{ref3};
$\triangle$, ref \citenum{ref5}; $+$, ref \citenum{ref43}. All symbols are BET
model values evaluated for the solubility data of the reference literature.}
\end{figure}

\begin{table}[htbp]
\centering
\caption{Deviation of the Vapor Pressures of the LiCl$\cdot$H$_2$O Saturated
Solution Calculated for Different Solubility Data from the Standard Values}
\begin{tabular}{lccc}
\toprule
literature of solubility data
& \begin{tabular}[c]{@{}c@{}}temperature\\ range, $T$/K\end{tabular}
& \begin{tabular}[c]{@{}c@{}}number of\\ data, $n$\end{tabular}
& \begin{tabular}[c]{@{}c@{}}average\\ deviation\textsuperscript{a}\end{tabular} \\
\midrule
Linke and Seidell\textsuperscript{1} & 298.15--369.15 & 9 & 0.96 \\
Farelo et al.\textsuperscript{43} & 296.7--334.5 & 10 & 0.97 \\
Clynne and Potter\textsuperscript{3} & 297.95--365.1 & 9 & 1.04 \\
Monnin et al.\textsuperscript{5} & 298.15--363.15 & 14 & 1.08 \\
Friend and Colley\textsuperscript{37} & 291.15--361.15 & 13 & 1.49 \\
Benrath\textsuperscript{35} & 298.15--363.15 & 5 & 1.72 \\
Johnson, Jr.\ and Molstad\textsuperscript{9} & 303.15, 323.15 & 2 & 2.38 \\
H\"uttig and Reuscher\textsuperscript{6} & 303.15--100.5 & 6 & 2.65 \\
International Critical Tables\textsuperscript{42} & 293.15--373.65 & 10 & 4.92 \\
\bottomrule
\end{tabular}

\vspace{2pt}
\begin{minipage}{0.8\textwidth}
\footnotesize \textsuperscript{a} Average deviation
$= \frac{1}{n} \sum_{i=1}^{n} |p_i^{\mathrm{calc}} - p_i^{\mathrm{std}}| /
p_i^{\mathrm{std}} \times 100$.
\end{minipage}
\end{table}

\begin{figure}[htbp]
\centering
% \includegraphics[width=0.6\textwidth]{figure13}
\caption{Vapor pressures of the LiCl$\cdot$H$_2$O saturated solution vs
temperature: ---, standard values; \ldots, model uncertainty based on the
standard values; $\square$, ref \citenum{ref6}, $\triangle$,
ref \citenum{ref9}; $\triangledown$, ref \citenum{ref42}; $\diamond$,
ref \citenum{ref35}; $+$, ref \citenum{ref37}. All symbols are BET model
values evaluated for the solubility data of the reference literature.}
\end{figure}

\subsection*{On the Solubility Data of LiCl}

The solubility data of the solid-phase LiCl given by different
authors\rc{ref1,ref3,ref5,ref6,ref11,ref35,ref42,ref44} do not agree with each
other, as shown in Figure 14. It is also noted that the vapor pressures
calculated for these solubility data are generally lower than the standard
curve (Figure 15 and Table 5). This result may suggest the uncertainty of the
model and/or the standard values of the vapor pressure. On one hand, the
standard vapor pressure data are fitted only to the experimental data of
Applebey et al.,\rc{ref11} the uncertainty range of which was not reported.
One the other hand, the BET model parameters at 428.65 K are fitted to the
experimental data of Brendler and Voigt,\rc{ref31} who experimentally
determined the isopiestic molalities of LiCl solution with a MgCl$_2$
reference solution and calculated the osmotic coefficients with a set of
Pitzer model parameters\rc{ref45} for the MgCl$_2$ solution. Simultaneously,
they also compared the model values with Dittrich's experimental
values,\rc{ref46} yielding an average deviation of $\pm0.04$ for the osmotic
coefficient. Thus, it can be reasonably assumed that the osmotic coefficient
of LiCl solution determined by the isopiestic method should also have an
uncertainty of $\pm0.04$. This uncertainty is drawn in Figure 15 on the basis
of the standard curve. Of all of the calculated vapor pressures for the
solubility data,\rc{ref1,ref3,ref5,ref6,ref11,ref35,ref42,ref44} only those
pertaining to the data of Clynne et al.\rc{ref3} and Benrath\rc{ref35} fall in
the uncertainty range of the standard curve. The exact differences between the
calculated and standard values are presented in the last column of Table 5.
Although Clynne et al.'s data\rc{ref3} are different from most of the others,
their data should not be discounted when one evaluates the solubility data of
LiCl critically. This conclusion can also find support from the visual
measurement method of Clynne et al.,\rc{ref3,ref47} who suggested that the
chemical analysis of the highly viscous saturated solutions gives larger
values because of suspended crystallites. Although Friend et al.\rc{ref37}
also employed the visual method, Clynne et al.'s data,\rc{ref3} with an
average precision of 0.09 \%, may be more accurate than Friend et al.'s data.

\begin{figure}[htbp]
\centering
% \includegraphics[width=0.6\textwidth]{figure14}
\caption{Solubility of the solid-phase LiCl(s) in pure water vs temperature:
$\circ$, ref \citenum{ref1}; $\blacktriangle$, ref \citenum{ref3};
$\triangle$, ref \citenum{ref5}; $\diamond$, ref \citenum{ref6}; $+$,
ref \citenum{ref11}; $\blacksquare$, ref \citenum{ref35}; $*$,
ref \citenum{ref42}; $\times$, ref \citenum{ref44}.}
\end{figure}

\begin{figure}[htbp]
\centering
% \includegraphics[width=0.6\textwidth]{figure15}
\caption{Vapor pressures of the LiCl saturated solution vs temperature:
---, standard values; \ldots, model uncertainty based on the standard values;
$\circ$, ref \citenum{ref1}; $\blacktriangle$, ref \citenum{ref3};
$\triangle$, ref \citenum{ref5}; $\diamond$, ref \citenum{ref6}; $+$,
ref \citenum{ref11}; $\blacksquare$, ref \citenum{ref35}; $*$,
ref \citenum{ref42}; $\times$, ref \citenum{ref44}. All symbols are BET model
values evaluated for the solubility data of the reference literature.}
\end{figure}

\begin{table}[htbp]
\centering
\caption{Deviation of the Vapor Pressures $p$ of the LiCl Saturated Solution
Calculated for Different Solubility Data from the Standard Values}
\begin{tabular}{lccc}
\toprule
literature of solubility data
& \begin{tabular}[c]{@{}c@{}}temperature\\ range, $T$/K\end{tabular}
& \begin{tabular}[c]{@{}c@{}}number of\\ data, $n$\end{tabular}
& \begin{tabular}[c]{@{}c@{}}average\\ deviation\textsuperscript{a}\end{tabular} \\
\midrule
Clynne and Potter\textsuperscript{3} & 385.65--428.6 & 5 & 2.718 \\
Benrath\textsuperscript{35} & 372.65 & 1 & 1.675 \\
Applebey et al.\textsuperscript{11} & 373.35--428.75 & 4 & 4.577 \\
Monnin et al.\textsuperscript{5} & 373--433 & 14 & 5.672 \\
Friend et al.\textsuperscript{44} & 374.95--427.15 & 9 & 5.697 \\
Linke and Seidell\textsuperscript{1} & 369.15--433.15 & 7 & 5.810 \\
International Critical Tables\textsuperscript{42} & 373.65--393.15 & 3 & 5.945 \\
H\"uttig and Reuscher\textsuperscript{6} & 373.35--433.15 & 3 & 7.208 \\
\bottomrule
\end{tabular}

\vspace{2pt}
\begin{minipage}{0.8\textwidth}
\footnotesize \textsuperscript{a} Average deviation
$= \frac{1}{n} \sum_{i=1}^{n} |p_i^{\mathrm{calc}} - p_i^{\mathrm{std}}| /
p_i^{\mathrm{std}} \times 100$.
\end{minipage}
\end{table}

\section*{Conclusions}

By fitting the BET model parameters to the selected water activity data in the
system LiCl + H$_2$O at salt concentrations below 20 mol$\cdot$kg$^{-1}$ and
in the temperature range from 273.15 K to 428.65 K, the vapor pressures for
solubility data given by different authors were calculated and compared with
reliable experimental data. By comparing the consistency of calculated and
experimental vapor pressures above LiCl$\cdot n$H$_2$O saturated solution, the
reliability of the solubility data was evaluated. The investigation of this
work indicated that the vapor pressures evaluated for the solubility data of
LiCl$\cdot$2H$_2$O given by Clynne et al., Monnin et al., Linke and Seidell,
except that at two Linke and Seidell points at 273.15 K and 278.15 K, are
generally in good agreement with the experimental data within the model
estimated error. For all other LiCl$\cdot$2H$_2$O solubility data cited in the
work, the calculated vapor pressures are scattered. A comparison of the
calculated and experimental vapor pressures also indicated that all of the
solubility data for LiCl$\cdot$H$_2$O from Linke and Seidell, Clynne and
Potter, and Monnin et al. are reliable. However, the calculated vapor
pressures above the LiCl saturated solution are systematically lower than the
experimental data. Among them, only the vapor pressures calculated using
Clynne and Potter's solubility data are consistent with the experimental
values in the error range of the model. Consequently, Clynne and Potter's
solubility data for LiCl should not be discounted, although their data are
different from most of the others' data.

\begin{thebibliography}{99}
\footnotesize

\bibitem{ref1} Linke, W. F.; Seidell, A. \textit{Solubilities: Inorganic and
Metal--Organic Compounds}, 4th ed.; American Chemical Society: Washington, DC,
1965.

\bibitem{ref2} Schimmel, F. A. Solubility of lithium chloride and lithium
thiocyanate at low temperatures. \textit{J. Chem. Eng. Data} \textbf{1960},
\textit{5}, 519--520.

\bibitem{ref3} Clynne, M. A.; Potter, R. W., II. Solubility of some alkali and
alkaline earth chlorides in water at moderate temperatures. \textit{J. Chem.
Eng. Data} \textbf{1979}, \textit{24}, 338--340.

\bibitem{ref4} Cohen-Adad, R.; Lorimer, J. W. \textit{Alkali Metal and
Ammonium Chlorides in Water and Heavy Water}; Pergamon: New York, 1991.

\bibitem{ref5} Monnin, C.; Dubois, M.; Papaiconomou, N.; Simonin, J.-P.
Thermodynamics of the LiCl--H$_2$O system. \textit{J. Chem. Eng. Data}
\textbf{2002}, \textit{47}, 1331--1336.

\bibitem{ref6} H\"uttig, G. H.; Reuscher, F. Study on the chemistry of
Lithium. I. On the hydration of LiCl and LiBr. \textit{Z. Anorg. Allg. Chem.}
\textbf{1924}, \textit{137}, 155--159.

\bibitem{ref7} Pearce, J. N.; Nelson, A. F. The vapour pressures of aqueous
solution of lithium nitrate and activity coefficients of some alkali salts in
solutions of high concentration at 25 $^{\circ}$C. \textit{J. Am. Chem. Soc.}
\textbf{1932}, \textit{54}, 3544--3555.

\bibitem{ref8} Rockland, L. B. Saturated salt solutions for static control of
relative humidity between 5 $^{\circ}$C and 40 $^{\circ}$C. \textit{Anal.
Chem.} \textbf{1960}, \textit{32}, 1375--1376.

\bibitem{ref9} Johnson, E. F., Jr.; Molstad, M. C. Thermodynamic properties of
aqueous lithium chloride solutions. \textit{J. Phys. Colloid Chem.}
\textbf{1951}, \textit{55}, 257--281.

\bibitem{ref10} \textit{Gmelings Handbuch Der Anorganischen Chemie}; Lithium,
supplemental band, system number 20; Verlag Chemie, GmbH: Weinheim, Germany,
1960.

\bibitem{ref11} Applebey, M. P.; Crawford, F. H.; Gordon, K. Vapour pressures
of saturated solutions. Lithium chloride and lithium sulphate. \textit{J.
Chem. Soc.} \textbf{1934}, 1665--1671.

\bibitem{ref12} Hedlin, C. P.; Trofimenkoff, F. N. In \textit{Humidity and
Moisture}; Wexler, A., Ed.; Reinhold Publishing: New York, 1965; Vol. 3,
pp 519--520.

\bibitem{ref13} Acheson, D. T. In \textit{Humidity and Moisture}; Wexler, A.,
Ed.; Reinhold Publishing: New York, 1965; Vol. 3, pp 521--530.

\bibitem{ref14} Stokes, R. H.; Robinson, R. A. Standard solutions for humidity
control at 25 $^{\circ}$C. \textit{Ind. Eng. Chem.} \textbf{1949},
\textit{41}, 3013.

\bibitem{ref15} Morillon, V.; Debeaufort, F.; Jose, J. Tharrault, J. F.;
Capelle, M.; Blond, G.; Voilley, A. Water vapour pressure above saturated salt
solutions at low temperatures. \textit{Fluid Phase Equilib.} \textbf{1999},
\textit{155}, 297--309.

\bibitem{ref16} Greenspan, L. Humidity fixed points of binary saturated
aqueous solutions, \textit{J. Res. Natl. Bur. Stand. (U.S.)} \textbf{1977},
\textit{81A}, 89--96.

\bibitem{ref17} Clegg, S. L.; Pitzer, K. S. Thermodynamics of multicomponent,
miscible, ionic solutions: generalized equations for systemetrical
electrolytes. \textit{J. Phys. Chem.} \textbf{1992}, \textit{96}, 3513--3520.

\bibitem{ref18} Clegg, S. L.; Pitzer, K. S.; Brimblecombe, P. Thermodynamics
of multicomponent, miscible, ionic solutions. 2. Mixtures including
unsymmetrical electrolytes. \textit{J. Phys. Chem.} \textbf{1992}, \textit{96},
9470--9479.

\bibitem{ref19} Pitzer, K. S.; Simonson, J. M. Thermodynamics of
multicomponent, miscible, ionic systems: theory and equations. \textit{J. Phys.
Chem.} \textbf{1986}, \textit{90}, 3005--3009.

\bibitem{ref20} Stokes, R. H.; Robinson, R. A. Ionic hydration and activity in
electrolyte solutions. \textit{J. Am. Chem. Soc.} \textbf{1948}, \textit{70},
1870--1878.

\bibitem{ref21} Sander, B.; Rasmussen, P.; Fredenslund, A. Calculation of
solid/liquid equilibria in aqueous solution of nitrate salts using an extended
uniquac equation. \textit{Chem. Eng. Sci.} \textbf{1986}, \textit{41},
1197--1202.

\bibitem{ref22} Chen, C.-C.; Mathias, P. M.; Orbey, H. Use of hydration and
dissociation chemistries with the electrolyte-NRTL model. \textit{AIChE J.}
\textbf{1999}, \textit{45}, 1576--1586.

\bibitem{ref23} Simonin, J.-P. Real ionic solution in the mean spherical
approximation. 2. Pure strong electrolyte up to very high concentrations, and
mixtures, in the primitive model. \textit{J. Phys. Chem. B} \textbf{1997},
\textit{101}, 4313--4320.

\bibitem{ref24} Gibbard, H. F., Jr.; Scatchard, G. Liquid--vapor equilibrium
of aqueous Lithium chloride, from 25 $^{\circ}$C to 100 $^{\circ}$C and from
1.0 to 18.5 mole, and related properties. \textit{J. Chem. Eng. Data}
\textbf{1973}, \textit{18}, 293--298.

\bibitem{ref25} Kangro, W.; Groeneveld, A. \textit{Z. Phys. Chem. N. F.}
\textbf{1962}, \textit{32}, 110--126.

\bibitem{ref26} Robinson, R. A.; Stokes, R. H. \textit{Electrolyte Solution},
2nd ed.; Butterworths: London, 1959.

\bibitem{ref27} Staples, B. R. Activity and osmotic coefficients of aqueous
sulphuric acid at 298.15 K. \textit{J. Chem. Ref. Data} \textbf{1981},
\textit{10}, 779--798.

\bibitem{ref28} Sako, T.; Hakuta, T.; Yoshitome, H. Vapor pressures of binary
(H$_2$O--HCl, --MgCl$_2$, and --CaCl$_2$) and ternary
(H$_2$O--MgCl$_2$--CaCl$_2$) aqueous solutions. \textit{J. Chem. Eng. Data}
\textbf{1985}, \textit{30}, 224--228.

\bibitem{ref29} Rard, J. A.; Wijesinghe, A. M.; Wolery, T. J. Review of the
thermodynamic properties of Mg(NO$_3$)$_2$(aq) and their representation with
the standard and extended ion-interaction (Pitzer) models at 298.15 K.
\textit{J. Chem. Eng. Data} \textbf{2004}, \textit{49}, 1127--1140.

\bibitem{ref30} Hill, P. G. A unified Fundamental equation for the
thermodynamic properties of H$_2$O. \textit{J. Phys. Chem. Ref. Data}
\textbf{1990}, \textit{19}, 1233--1274.

\bibitem{ref31} Brendler, V.; Voigt, W. Isopiestic measurements at high
temperatures. 1. Aqueous solutions of LiCl, CsCl, and CaCl$_2$ at
155 $^{\circ}$C. \textit{J. Solution Chem.} \textbf{1994}, \textit{23},
1061--1072.

\bibitem{ref32} Apelbat, A.; Korin, E. Vapour pressures of saturated aqueous
solutions of ammonium iodide, potassium iodide, potassium nitrate, strontium
chloride, lithium sulphate, sodium thiosulphate, magnesium nitrate, and uranyl
nitrate from $T = (278$ to 323) K. \textit{J. Chem. Thermodyn.} \textbf{1998},
\textit{30}, 459--471.

\bibitem{ref33} Stokes, R. H.; Robinson, R. A. Standard solutions for humidity
control at 25 $^{\circ}$C. \textit{Ind. Eng. Chem.} \textbf{1949},
\textit{41}, 2013.

\bibitem{ref34} Kessis, J. J. Syst\`eme eau-chlorure de lithium. Mesures de
solubilit\'e entre $-65$ et $+60$ $^{\circ}$C. \textit{Bull. Soc. Chim. Fr.}
\textbf{1961}, 1503--1504.

\bibitem{ref35} Benrath, H. Die polythermen der tern\"aren systeme:
CuCl$_2$--(LiCl)$_2$--H$_2$O und NiCl$_2$--(LiCl)$_2$--H$_2$O. \textit{Z.
Anorg. Allg. Chem.} \textbf{1932}, \textit{205}, 417--424.

\bibitem{ref36} Moran, H. E., Jr. System lithium chloride--water. \textit{J.
Phys. Chem.} \textbf{1956}, \textit{60}, 1666--1667.

\bibitem{ref37} Friend, J. A. N.; Colley, A. T. W. The solubility of lithium
chloride in water. \textit{J. Chem. Soc.} \textbf{1931}, \textit{31},
48--3149.

\bibitem{ref38} Bassett, H.; Sanderson, I. The compounds of lithium chloride
with cobalt chloride. Water as a linking agent in polynuclear cations.
\textit{J. Chem. Soc.} \textbf{1932}, 1855--1864.

\bibitem{ref39} Woskresenskaya, N. K.; Yanatieva, O. K. \textit{Izv. Akad.
Nauk SSSR, Ser. Khim.} \textbf{1937}, 97--121.

\bibitem{ref40} Novoselova, A. V.; Sosnovskaya, I. G. \textit{Zh. Obshchei.
Khim.} \textbf{1951}, \textit{21}, 813--817.

\bibitem{ref41} Benrath, H. Das system
kobaltchlorid-lithiumchlorid-wasser. \textit{Z. Anorg. Allg. Chem.}
\textbf{1939}, \textit{240}, 87--96.

\bibitem{ref42} \textit{International Critical Tables}, 1st ed.; McGraw-Hill:
New York, 1928; Vol. 3, p 368.

\bibitem{ref43} Farelo, F.; Fernandes, C.; Avelino, A. Solubilities for six
ternary systems: NaCl + NH$_4$Cl + H$_2$O, KCl + NH$_4$Cl + H$_2$O,
NaCl + LiCl + H$_2$O, KCl + LiCl + H$_2$O, NaCl + AlCl$_3$ + H$_2$O, and
KCl + AlCl$_3$ + H$_2$O at $T = (298$ to 333) K. \textit{J. Chem. Eng. Data}
\textbf{2005}, \textit{50}, 1470--1477.

\bibitem{ref44} Friend, J. N.; Hale, R. W.; Ryder, E. A. The solubility of
lithium chloride in water between 70$^{\circ}$ and 160$^{\circ}$. \textit{J.
Chem. Soc.} \textbf{1937}, 970.

\bibitem{ref45} Valyashko, V. M.; Urosova, M. A.; Voigt, W.; Emons, H.-H.
Properties of the system MgCl$_2$--H$_2$O in a wide range of temperature and
concentration. \textit{Zh. Neorg. Khim.} \textbf{1988}, \textit{33},
228--232.

\bibitem{ref46} Dittrich, A. Ph.D. Thesis, Bergakademie Freiberg, 1986.

\bibitem{ref47} \textit{CODATA Thermodynamic Tables: Selections for Some
Compounds of Calcium and Related Mixtures: A Prototype Set of Tables}; Gavin,
D., Parker, V. B., White, H. J., Jr., Eds.; Hemisphere Publishing Corporation:
Washington, DC, 1987.

\end{thebibliography}

\end{document}
